\chapter*{Glossary of Terms and Abbreviations}
\addcontentsline{toc}{chapter}{Glossary of Terms and Abbreviations}

The following definitions cover all technical, statistical, and platform-specific
terms used in this report. They are written for readers without a statistical or
digital marketing background.

\vspace{8pt}
\begin{longtable}{@{}L{4cm} L{11.5cm}@{}}
\toprule
\textbf{Term / Acronym} & \textbf{Plain-English Definition} \\
\midrule
\endfirsthead
\multicolumn{2}{c}{\tablename\ \thetable{} (continued)} \\
\toprule
\textbf{Term / Acronym} & \textbf{Plain-English Definition} \\
\midrule
\endhead
\bottomrule
\endfoot
\bottomrule
\endlastfoot

\textbf{API} &
Application Programming Interface. A connection that lets software retrieve
data from another service automatically. Here it refers to the Apify API, which
is used to collect public Instagram data. \\[4pt]

\textbf{Apify} &
A web data collection service. This project used Apify's Instagram Scraper to
gather public post, reel, comment, and mention data for the account analysed in
this report. \\[4pt]

\textbf{Bootstrap / Bootstrap CI} &
A statistical method that estimates uncertainty by repeatedly re-drawing random
samples from the data and measuring how much the result varies. A \textbf{bootstrap
confidence interval (CI)} is a range of values that is likely to contain the true
answer 95\% of the time. \\[4pt]

\textbf{Call to Action (CTA)} &
A direct instruction in a caption or story that tells the viewer what to do next.
Examples: ``Book now'', ``Watch the full piece'', ``Link in bio to get involved''. \\[4pt]

\textbf{Coefficient of Variation (CV)} &
A measure of how consistent a group of numbers is. A CV of 0.30 means the
standard deviation is 30\% of the mean. A high CV (above 0.5) means results are
unpredictable and driven by occasional spikes. \\[4pt]

\textbf{Cohen's \textit{d}} &
A measure of \textit{practical} effect size. It answers: even if the difference is
statistically real, is it big enough to matter? Values: below 0.2 = negligible,
0.2--0.5 = small, 0.5--0.8 = medium, above 0.8 = large. \\[4pt]

\textbf{Confidence Interval (CI)} &
A range around an estimate that expresses uncertainty. A 95\% CI means: if the
analysis were repeated many times on new data, the true value would fall inside
this range 95\% of those times. Wider intervals mean less certainty; narrower
intervals mean more. \\[4pt]

\textbf{Content Pillar} &
A thematic category used to group posts by topic, derived from each account's own
content (for example projects, behind-the-scenes, events, or brand posts). Pillar
analysis shows which topics attract the most audience interaction. \\[4pt]

\textbf{Degrees of Freedom (df)} &
A technical input to statistical tests. In plain terms, it reflects the amount
of information available. A higher df generally means a more reliable test result.
It is shown in test tables for completeness but does not require interpretation
by the reader. \\[4pt]

\textbf{Direct Message (DM)} &
A private message sent between Instagram users. DMs are not visible publicly and
are not captured in this analysis. They are likely a significant enquiry and
conversion channel that this report cannot measure. \\[4pt]

\textbf{Engagement Score} &
A composite score used throughout this report to compare posts on a common scale.
It is calculated as: $\text{Likes} + (\text{Comments} \times 5) + (\text{Video Plays} \div 100)$.
Comments are weighted more heavily because they require active effort and carry
stronger commercial signal than a passive like or view. \\[4pt]

\textbf{Exploratory Data Analysis (EDA)} &
An initial examination of a dataset to understand its shape, range, and
characteristics before running formal tests. In Part II of this report, EDA
includes summary tables covering mean, median, standard deviation, and skewness
for each content segment. \\[4pt]

\textbf{FAQ} &
Frequently Asked Questions. In the context of Instagram Highlights, an FAQ
Highlight is a saved Story that answers common audience questions (location,
hours, pricing, how to book or get involved) so that a viewer can find the answer
without having to comment or send a message. \\[4pt]

\textbf{Hypothesis Test} &
A formal statistical procedure for deciding whether an observed difference is
likely to be real or likely to be the result of random chance. Each test
produces a \textit{p}-value (see below). In this report, hypothesis tests are
used to evaluate claims such as ``short videos attract more engagement than
feed posts''. \\[4pt]

\textbf{Interquartile Range (IQR)} &
The range between the 25th and 75th percentile of a dataset. It captures the
middle 50\% of values and is used in outlier detection: a post is flagged as an
outlier if its engagement score exceeds the 75th percentile plus 1.5 times the IQR. \\[4pt]

\textbf{Kruskal-Wallis Test} &
A non-parametric statistical test used to compare three or more groups. It does
not assume the data is normally distributed, making it appropriate for social
media engagement scores, which are heavily skewed. In this report it is used to
test whether posting day has a real effect on engagement. \\[4pt]

\textbf{KPI} &
Key Performance Indicator. A metric chosen to track progress toward a business
goal. This report recommends specific KPIs such as public engagement score by
pillar, comment volume, and response time to commercial comments. \\[4pt]

\textbf{Mann-Whitney U Test} &
A non-parametric test that compares two groups by ranking all observations
together and checking whether one group tends to rank higher than the other.
It does not require normal distributions and is the preferred test when data
is skewed, as social media engagement data typically is. \\[4pt]

\textbf{Mean} &
The arithmetic average of a set of values. It is sensitive to extreme outliers.
In this report, means are always shown alongside medians so that the effect of
outlier posts can be seen. \\[4pt]

\textbf{Median} &
The middle value when all observations are sorted from lowest to highest.
It is more robust than the mean when data is skewed. If the median engagement
score is, say, 40, that means half of all posts scored below 40 and half scored
above. \\[4pt]

\textbf{Monte Carlo Simulation (MC)} &
A method of estimating outcomes by running thousands of random trials drawn from
real historical data, then summarising the distribution of results. It is used
here to project expected engagement under different content strategies (MC1),
forecast engagement over 30, 60, and 90 days (MC2), model the booking conversion
pipeline (MC3), optimise pillar allocations (MC4), and calculate the probability
of achieving specific targets (MC5). \\[4pt]

\textbf{$n$} &
The number of observations in a sample. In the context of a pillar analysis,
$n = 8$ for a pillar means only eight posts in the dataset were classified under
that pillar. Small $n$ leads to wide confidence intervals and unreliable
conclusions. \\[4pt]

\textbf{$p$-value} &
The probability that the observed result (or one more extreme) would occur by
chance alone, assuming there is no real effect. A $p$-value below 0.05 is the
conventional threshold for declaring a result ``statistically significant''.
A $p$-value of 0.07, for example, means the result is not conclusive but is
directionally informative. \\[4pt]

\textbf{Pearson Correlation ($r$)} &
A measure of the straight-line relationship between two variables, ranging from
$-1$ (perfect negative) to $+1$ (perfect positive). A value near zero means the
two variables are not meaningfully related. For example, a Pearson $r$ near 0.10
between caption length and engagement would indicate that longer captions do not
reliably produce more engagement. \\[4pt]

\textbf{Percentile (P5, P25, P75, P95)} &
A value below which a given percentage of observations fall. P5 means the bottom
5\% of results; P95 means the top 5\%. In the Monte Carlo tables, P5 represents
the worst-case outcome seen in 5\% of simulations and P95 represents the best-case
outcome seen in 5\% of simulations. \\[4pt]

\textbf{Reel} &
Instagram's short-video format, typically 15 to 90 seconds. Reels are distributed
by Instagram's recommendation algorithm and therefore reach audiences beyond the
account's existing followers. Feed posts (images and carousels) are distributed
primarily to existing followers. \\[4pt]

\textbf{Seed (Random Seed)} &
A fixed starting number given to a random number generator so that the same
sequence of random values is produced every time the code runs. Setting the seed
to 20260530 means the Monte Carlo results in this report are exactly reproducible. \\[4pt]

\textbf{Service Level Agreement (SLA)} &
A defined standard of service, here applied to response time. An Instagram comment
SLA of 30 minutes means the team commits to replying to all commercial comments
within 30 minutes of a comment being made. \\[4pt]

\textbf{Shortcode} &
The unique string in an Instagram post URL. For the URL
\texttt{https://www.instagram.com/p/DN7\_M4B2HMT/}, the shortcode is
\texttt{DN7\_M4B2HMT}. It is used in this analysis as the unique identifier for
each post when deduplicating the dataset. \\[4pt]

\textbf{Skewness} &
A measure of how asymmetric a distribution is. A positive skew (the typical
pattern for social media engagement) means that most posts score low, with a small
number of high-performing posts pulling the mean upward. \\[4pt]

\textbf{Spearman Correlation ($\rho$, rho)} &
A version of the Pearson correlation that works with ranked values rather than raw
numbers. It is more robust to outliers and skewed data. In this report it is used
alongside Pearson $r$ to confirm directional findings. \\[4pt]

\textbf{Standard Deviation (Std Dev)} &
A measure of how spread out values are around the mean. A high standard deviation
means results are inconsistent; a low standard deviation means they are
predictable. \\[4pt]

\textbf{Statistically Significant} &
A result is statistically significant when the $p$-value is below the chosen
threshold (0.05 in this report). It means the observed effect is unlikely to have
occurred by chance alone. It does not necessarily mean the effect is large or
commercially important. \\[4pt]

\textbf{UGC (User-Generated Content)} &
Content created by customers, fans, or third parties rather than the account
owner. Third-party mentions of the account are a form of UGC. Reposting UGC with
attribution adds social proof without requiring original production effort. \\[4pt]

\textbf{UTC (Coordinated Universal Time)} &
The global time standard. All post timing in this report uses UTC. For accounts
based in GMT/UTC+0 regions (such as Ghana, which observes GMT year-round), the
hour values in the timing tables correspond directly to local time; accounts in
other time zones should offset accordingly. \\[4pt]

\textbf{Welch t-test} &
A parametric statistical test that compares the means of two groups. Unlike the
standard t-test, it does not assume the two groups have equal variance, making it
more appropriate for comparing groups of different sizes. \\[4pt]

\textbf{Bonferroni Correction} &
An adjustment made when running multiple hypothesis tests at the same time. It
lowers the $p$-value threshold (from 0.05 to 0.0083 when running six tests) to
reduce the chance that at least one result appears significant purely by accident. \\[4pt]

\textbf{Excess Kurtosis} &
A measure of how ``heavy-tailed'' a distribution is compared to a normal bell
curve. Positive excess kurtosis means extreme values (very high-performing posts)
occur more often than a normal distribution would predict. A large positive value
indicates the distribution is dominated by a few outlier posts. \\[4pt]

\textbf{Lift vs Baseline} &
The percentage difference between a pillar's average engagement score and the
overall average. A lift of +100\% for a pillar means that, on average, a post in
that pillar attracted twice the engagement of a typical post across all
categories. \\[4pt]

\textbf{P(achieve)} &
In the Monte Carlo risk tables, this is the probability (expressed as a
percentage) that a strategy will reach a specified engagement target over 12 weeks.
P(achieve) of 80\% means that in 8 out of every 10 simulated runs, the strategy
reached the target. \\[4pt]

\textbf{Top-20 Hit Rate} &
The proportion of posts within a pillar that scored in the top 20\% of all posts
by engagement score. A high hit rate means the pillar consistently produces
strong content; a low hit rate means performance is variable. \\[4pt]

\end{longtable}
